% ============================================================
% 附录 (Appendices)
% ============================================================
\clearpage
\phantomsection
\addcontentsline{toc}{chapter}{附录}
\chapter*{附录}
\markboth{附录}{附录}

% ============================================================
% 附录 A：数据完整备查表
% ============================================================
\section*{附录 A：数据完整备查表}
\addcontentsline{toc}{section}{附录 A：数据完整备查表}

本附录提供第 4 章正文中"精简版"表格所对应的完整数据表。

% ---------- 附表 A-1：基线特征完整表 ----------
\begin{landscape}
\thispagestyle{empty}
\begin{table}[htbp]
\centering
\caption{研究二 PP 样本基线特征完整比较表（n = 24）}
\label{tab:app_baseline_full}
\zihao{5}
\begin{threeparttable}
\begin{tabular}{lccccccc}
\toprule
变量 & AI组 & Self组 & 均值差 & 95\% CI & t检验 p & MWU p & Hedges' g \\
\midrule
年龄（岁）                & 24.91 $\pm$ 2.55     & 23.85 $\pm$ 1.63     & +1.06  & (--0.82, 2.95)    & 0.250 & 0.165 & 0.49 \\
身高（cm）                & 179.73 $\pm$ 7.13    & 176.00 $\pm$ 4.58    & +3.73  & (--1.55, 9.01)    & 0.154 & 0.123 & 0.61 \\
体重（kg）                & 77.54 $\pm$ 10.20    & 74.27 $\pm$ 4.79     & +3.27  & (--3.93, 10.47)   & 0.346 & 0.361 & 0.41 \\
BMI（kg/m$^2$）           & 23.91 $\pm$ 1.73     & 23.97 $\pm$ 1.05     & --0.06 & (--1.33, 1.21)    & 0.926 & 0.531 & --0.04 \\
体脂率（\%）              & 15.49 $\pm$ 2.19     & 15.53 $\pm$ 2.09     & --0.04 & (--1.87, 1.79)    & 0.964 & 0.954 & --0.02 \\
骨骼肌量（kg）            & 37.16 $\pm$ 4.79     & 35.15 $\pm$ 2.35     & +2.02  & (--1.38, 5.41)    & 0.224 & 0.173 & 0.53 \\
去脂体重（kg）            & 65.47 $\pm$ 8.36     & 62.71 $\pm$ 3.88     & +2.76  & (--3.13, 8.66)    & 0.331 & 0.303 & 0.42 \\
基线深蹲 1RM（kg）        & 130.68 $\pm$ 12.15   & 130.77 $\pm$ 17.45   & --0.09 & (--12.70, 12.53)  & 0.989 & 0.662 & --0.01 \\
基线相对 1RM（kg/kg）     & 1.70 $\pm$ 0.19      & 1.76 $\pm$ 0.23      & --0.06 & (--0.24, 0.12)    & 0.492 & 0.955 & --0.27 \\
基线 CMJ（cm）            & 59.66 $\pm$ 5.23     & 60.52 $\pm$ 6.27     & --0.85 & (--5.72, 4.02)    & 0.720 & 0.820 & --0.14 \\
基线 SJ（cm）             & 55.68 $\pm$ 4.89     & 54.81 $\pm$ 4.82     & +0.87  & (--3.26, 5.01)    & 0.665 & 0.569 & 0.17 \\
基线训练自我效能（分）    & 65.45 $\pm$ 5.09     & 65.69 $\pm$ 6.91     & --0.24 & (--5.33, 4.86)    & 0.924 & 0.977 & --0.04 \\
\bottomrule
\end{tabular}
\begin{tablenotes}[flushleft]\footnotesize
\item 注：连续变量以均值 $\pm$ 标准差表示；组间比较采用 Welch 独立样本 t 检验（主检验）与 Mann-Whitney U 检验（敏感性检验）并列报告。
\end{tablenotes}
\end{threeparttable}
\end{table}
\end{landscape}

% ---------- 附表 A-2：ANCOVA 五项结局完整表 ----------
\begin{landscape}
\thispagestyle{empty}
\begin{table}[htbp]
\centering
\caption{探索性结局指标 ANCOVA 结果完整表（n = 24）}
\label{tab:app_ancova_full}
\zihao{5}
\begin{threeparttable}
\begin{tabular}{lccccccc}
\toprule
结局指标 & AI组T1 & Self组T1 & 调整后差值 & 差值 95\% CI & p 值 & Hedges' g（95\% CI） & FDR校正 p \\
\midrule
深蹲绝对 1RM（kg）        & 142.27 $\pm$ 17.48 & 139.81 $\pm$ 18.18 & +2.43  & (--2.01, 6.87)   & 0.266     & 0.46 (--0.35, 1.28) & 0.341 \\
深蹲相对 1RM（kg/kg）     & 1.85 $\pm$ 0.24    & 1.88 $\pm$ 0.24    & +0.03  & (--0.03, 0.09)   & 0.273     & 0.39 (--0.42, 1.20) & 0.341 \\
CMJ 高度（cm）            & 65.06 $\pm$ 6.53   & 63.70 $\pm$ 6.21   & +2.16  & (--0.01, 4.33)   & 0.051     & 0.85 (0.02, 1.69)   & 0.126 \\
SJ 高度（cm）             & 59.98 $\pm$ 6.02   & 58.18 $\pm$ 5.46   & +0.39  & (--1.24, 2.02)   & 0.622     & 0.43 (--0.38, 1.24) & 0.622 \\
训练自我效能（分） & 83.18 $\pm$ 6.32 & 73.69 $\pm$ 7.15 & +9.88  & (6.86, 12.90) & $<$0.001** & 2.70 (1.59, 3.80) & $<$0.001** \\
\bottomrule
\end{tabular}
\begin{tablenotes}[flushleft]\footnotesize
\item 注：** p $<$ 0.001。ANCOVA 模型设定为 Post $\sim$ Group $+$ Pre $+$ Stratum，采用 Type III 偏差平方和。
\end{tablenotes}
\end{threeparttable}
\end{table}
\end{landscape}

% ---------- 附表 A-3：访谈受试者信息完整表 ----------
\begin{landscape}
\thispagestyle{empty}
\begin{table}[htbp]
\centering
\caption{质性访谈受试者基本信息与抽样特征完整分布表（n = 12）}
\label{tab:app_interview_full}
\zihao{5}
\begin{threeparttable}
\begin{tabular}{lcccccc}
\toprule
编号 & 组别 & 分层 & Delta1RM & DeltaCMJ & 训练执行特征 & 类型 \\
\midrule
P004 & AI 辅助组   & 低力量层 & +15.0 kg & +6.7 cm  & 命中率 85.0\%，未触发减组          & 积极接受型 \\
P010 & AI 辅助组   & 低力量层 & +12.5 kg & +3.1 cm  & 触发自动减组 1 次（S1）           & 安全依赖型 \\
P016 & AI 辅助组   & 低力量层 & +5.0 kg  & --0.3 cm & 触发自动减组 2 次（S4 与 S6）      & 安全依赖型 \\
P018 & AI 辅助组   & 低力量层 & +7.5 kg  & +10.3 cm & 命中率 74.2\%，SUS 偏低            & 批判接受型 \\
P025 & AI 辅助组   & 低力量层 & +22.5 kg & +6.7 cm  & 命中率 89.5\%，未触发减组          & 积极接受型 \\
P003 & 自我指导组  & 低力量层 & +7.5 kg  & +3.8 cm  & 自主调整 2 次，保守选重            & 规则内化型 \\
P012 & 自我指导组  & 高力量层 & +10.0 kg & +0.7 cm  & 经验丰富，依靠关节感觉调节          & 批判接受型 \\
P013 & 自我指导组  & 低力量层 & +10.0 kg & +8.1 cm  & 自主调整 4 次，严格维持动作质量    & 规则内化型 \\
P014 & 自我指导组  & 高力量层 & +10.0 kg & +1.5 cm  & 训练年限长，保守避免力竭          & 批判接受型 \\
P015 & 自我指导组  & 高力量层 & +17.5 kg & +0.6 cm  & 倾向于接近力竭与慢速硬顶          & 负荷导向型 \\
P021 & 自我指导组  & 低力量层 & +10.0 kg & +0.1 cm  & RIR 估计存在波动                  & 负荷导向型 \\
P027 & 自我指导组  & 低力量层 & +7.5 kg  & +3.6 cm  & 自主调整 3 次，依疲劳感减量      & 规则内化型 \\
\bottomrule
\end{tabular}
\begin{tablenotes}[flushleft]\footnotesize
\item 注：访谈对象采用最大差异抽样与主题饱和相结合的策略确定。
\end{tablenotes}
\end{threeparttable}
\end{table}
\end{landscape}

% ============================================================
% 附录 B：原始数据备查表
% ============================================================
\clearpage
\section*{附录 B：原始数据备查表}
\addcontentsline{toc}{section}{附录 B：原始数据备查表}

% ---------- 附表 B-1：SUS 与 TAM 个人评分明细 ----------
\begin{table}[htbp]
\centering
\caption{AI 辅助组系统评价量表描述性统计与个人评分明细（n = 11）}
\label{tab:app_sus_full}
\zihao{5}
\begin{threeparttable}
\begin{tabular}{lcccc}
\toprule
受试者编号 & SUS 总分 & PU 均分 & Trust 均分 & Intention 均分 \\
\midrule
P001               & 67.5 & 4.50 & 4.50 & 4.67 \\
P004               & 72.5 & 4.75 & 4.50 & 4.67 \\
P005               & 57.5 & 4.25 & 4.25 & 4.33 \\
P006               & 67.5 & 4.50 & 4.75 & 4.67 \\
P007               & 65.0 & 4.50 & 4.50 & 4.33 \\
P009               & 62.5 & 4.25 & 4.25 & 4.33 \\
P010               & 67.5 & 4.25 & 4.25 & 4.33 \\
P011               & 62.5 & 4.50 & 4.50 & 4.67 \\
P016               & 67.5 & 4.50 & 4.50 & 4.33 \\
P018               & 57.5 & 4.00 & 4.00 & 4.00 \\
P025               & 67.5 & 4.50 & 4.50 & 5.00 \\
\midrule
均值 $\pm$ 标准差 & 65.00 $\pm$ 4.61 & 4.41 $\pm$ 0.20 & 4.41 $\pm$ 0.20 & 4.49 $\pm$ 0.27 \\
范围             & 57.5--72.5 & 4.00--4.75 & 4.00--4.75 & 4.00--5.00 \\
\bottomrule
\end{tabular}
\begin{tablenotes}[flushleft]\footnotesize
\item 注：SUS = System Usability Scale（满分 100 分，预设推进阈值 $\ge$ 70 分）。
\end{tablenotes}
\end{threeparttable}
\end{table}

% ---------- 附表 B-2：28 对 Spearman 相关完整矩阵 ----------
\begin{table}[htbp]
\centering
\caption{训练负荷与结局变化的探索性 Spearman 相关完整矩阵（28 对组合）}
\label{tab:app_spearman_full}
\zihao{5}
\begin{threeparttable}
\begin{tabular}{llccc}
\toprule
训练执行变量 & 结局变量 & Spearman rho & 原始 p 值 & FDR 校正 p \\
\midrule
全期总负荷        & Delta绝对 1RM & +0.441 & 0.031* & 0.108 \\
深蹲主项总负荷    & Delta绝对 1RM & +0.398 & 0.054  & 0.126 \\
平均每课次负荷    & Delta绝对 1RM & +0.500 & 0.013* & 0.093 \\
总组数            & Delta绝对 1RM & --0.152 & 0.478  & 0.558 \\
总次数            & Delta绝对 1RM & --0.221 & 0.298  & 0.417 \\
训练总时长        & Delta绝对 1RM & +0.533 & 0.007* & 0.093 \\
个体 Hooper 均值  & Delta绝对 1RM & --0.455 & 0.025* & 0.102 \\
\midrule
全期总负荷        & Delta相对 1RM & +0.312 & 0.138 & 0.257 \\
深蹲主项总负荷    & Delta相对 1RM & +0.267 & 0.207 & 0.322 \\
平均每课次负荷    & Delta相对 1RM & +0.354 & 0.089 & 0.166 \\
总组数            & Delta相对 1RM & --0.098 & 0.649 & 0.706 \\
总次数            & Delta相对 1RM & --0.176 & 0.410 & 0.512 \\
训练总时长        & Delta相对 1RM & +0.398 & 0.054 & 0.126 \\
个体 Hooper 均值  & Delta相对 1RM & --0.331 & 0.114 & 0.220 \\
\midrule
全期总负荷        & DeltaCMJ & +0.212 & 0.321 & 0.437 \\
深蹲主项总负荷    & DeltaCMJ & +0.189 & 0.377 & 0.489 \\
平均每课次负荷    & DeltaCMJ & +0.245 & 0.249 & 0.373 \\
总组数            & DeltaCMJ & --0.052 & 0.809 & 0.821 \\
总次数            & DeltaCMJ & +0.108 & 0.615 & 0.706 \\
训练总时长        & DeltaCMJ & +0.267 & 0.207 & 0.322 \\
全期平均 sRPE     & DeltaCMJ & --0.198 & 0.354 & 0.472 \\
\midrule
全期总负荷        & DeltaSJ & +0.365 & 0.079 & 0.166 \\
深蹲主项总负荷    & DeltaSJ & +0.328 & 0.118 & 0.220 \\
平均每课次负荷    & DeltaSJ & +0.221 & 0.298 & 0.417 \\
总组数            & DeltaSJ & +0.479 & 0.018* & 0.093 \\
总次数            & DeltaSJ & +0.503 & 0.012* & 0.093 \\
训练总时长        & DeltaSJ & +0.494 & 0.014* & 0.093 \\
全期平均 sRPE     & DeltaSJ & --0.472 & 0.020* & 0.093 \\
\bottomrule
\end{tabular}
\begin{tablenotes}[flushleft]\footnotesize
\item 注：* 原始 p $<$ 0.05。FDR 校正采用 Benjamini-Hochberg 法（28 项检验共同校正）。
\end{tablenotes}
\end{threeparttable}
\end{table}

\clearpage
